\section{Ejercicios}

\subsection{Nivel básico}

\begin{ejercicio}

Una carga de prueba positiva de \(\SI{3.00}{\micro\coulomb}\) experimenta una fuerza eléctrica de \(\SI{0.0120}{\newton}\) hacia la derecha.

Determine la magnitud y dirección del campo eléctrico en ese punto.

\end{ejercicio}

\begin{respuesta}

\[
\boxed{E=\SI{4.00e3}{\newton\per\coulomb}}
\]

El campo apunta hacia la derecha.

\end{respuesta}

\begin{explicacion}

La definición de campo eléctrico es

\[
E=\frac{F}{q_0}.
\]

Convertimos la carga:

\[
q_0=\SI{3.00e-6}{\coulomb}.
\]

Sustituyendo,

\[
E=\frac{0.0120}{3.00\times10^{-6}}
=\SI{4.00e3}{\newton\per\coulomb}.
\]

Como la carga de prueba es positiva, el campo tiene el mismo sentido que la fuerza.

Por tanto,

\[
\boxed{E=\SI{4.00e3}{\newton\per\coulomb}}
\]

hacia la derecha.

\end{explicacion}


\begin{ejercicio}

Una carga

\[
q=+\SI{2.50}{\nano\coulomb}
\]

se coloca en una región donde existe un campo eléctrico uniforme de magnitud

\[
E=\SI{4.00e4}{\newton\per\coulomb}
\]

dirigido hacia la derecha.

Determine la fuerza eléctrica sobre la carga.

\end{ejercicio}

\begin{respuesta}

\[
\boxed{F=\SI{1.00e-4}{\newton}}
\]

hacia la derecha.

\end{respuesta}

\begin{explicacion}

La fuerza eléctrica sobre una carga situada en un campo es

\[
F=|q|E.
\]

Convertimos la carga:

\[
q=\SI{2.50e-9}{\coulomb}.
\]

Sustituyendo,

\[
F=(2.50\times10^{-9})(4.00\times10^4)
=\SI{1.00e-4}{\newton}.
\]

Como la carga es positiva, la fuerza tiene el mismo sentido que el campo.

\[
\boxed{F=\SI{1.00e-4}{\newton}}
\]

hacia la derecha.

\end{explicacion}


\begin{ejercicio}

Una carga

\[
q=-\SI{6.00}{\nano\coulomb}
\]

se encuentra en un campo eléctrico de magnitud

\[
E=\SI{2.50e5}{\newton\per\coulomb}
\]

dirigido hacia arriba.

Determine la fuerza eléctrica sobre la carga.

\end{ejercicio}

\begin{respuesta}

\[
\boxed{F=\SI{1.50e-3}{\newton}}
\]

hacia abajo.

\end{respuesta}

\begin{explicacion}

La magnitud de la fuerza es

\[
F=|q|E.
\]

Entonces,

\[
F=(6.00\times10^{-9})(2.50\times10^5)
=\SI{1.50e-3}{\newton}.
\]

Como

\[
q<0,
\]

la fuerza tiene sentido opuesto al campo.

Si el campo apunta hacia arriba, la fuerza sobre la carga negativa apunta hacia abajo.

Por tanto,

\[
\boxed{F=\SI{1.50e-3}{\newton}}
\]

hacia abajo.

\end{explicacion}


\begin{ejercicio}

Una carga puntual

\[
Q=+\SI{5.00}{\micro\coulomb}
\]

produce un campo eléctrico.

Determine la magnitud del campo a una distancia de \(\SI{0.200}{\metre}\).

\end{ejercicio}

\begin{respuesta}

\[
\boxed{E\approx\SI{1.12e6}{\newton\per\coulomb}}
\]

El campo apunta alejándose de la carga.

\end{respuesta}

\begin{explicacion}

Para una carga puntual,

\[
E=k\frac{|Q|}{r^2}.
\]

Sustituyendo,

\[
E=(8.99\times10^9)\frac{5.00\times10^{-6}}{(0.200)^2}.
\]

Entonces,

\[
\boxed{E\approx\SI{1.12e6}{\newton\per\coulomb}}.
\]

Como \(Q\) es positiva, el campo apunta radialmente hacia afuera.

\end{explicacion}


\begin{ejercicio}

Una carga puntual

\[
Q=-\SI{3.00}{\micro\coulomb}
\]

se encuentra en el origen.

Determine el campo eléctrico en un punto situado sobre el eje \(x\), a \(\SI{0.150}{\metre}\) a la derecha de la carga.

\end{ejercicio}

\begin{respuesta}

\[
\boxed{E\approx\SI{1.20e6}{\newton\per\coulomb}}
\]

hacia la izquierda.

\end{respuesta}

\begin{explicacion}

La magnitud del campo es

\[
E=k\frac{|Q|}{r^2}.
\]

Sustituyendo,

\[
E=(8.99\times10^9)\frac{3.00\times10^{-6}}{(0.150)^2}
\approx\SI{1.20e6}{\newton\per\coulomb}.
\]

Como la carga fuente es negativa, el campo apunta hacia ella.

El punto está a la derecha de la carga, por lo que el campo apunta hacia la izquierda.

\[
\boxed{E\approx\SI{1.20e6}{\newton\per\coulomb}}.
\]

\end{explicacion}


\begin{ejercicio}

Una carga puntual produce un campo eléctrico de magnitud \(E\) a cierta distancia \(r\).

Determine cómo cambia el campo si:

\begin{enumerate}
    \item la distancia se duplica;
    \item la distancia se triplica;
    \item la distancia se reduce a la mitad.
\end{enumerate}

\end{ejercicio}

\begin{respuesta}

\[
\boxed{r\to2r\Rightarrow E'=\frac{E}{4}}
\]

\[
\boxed{r\to3r\Rightarrow E'=\frac{E}{9}}
\]

\[
\boxed{r\to\frac{r}{2}\Rightarrow E'=4E}
\]

\end{respuesta}

\begin{explicacion}

El campo de una carga puntual satisface

\[
E\propto\frac{1}{r^2}.
\]

Si la distancia se multiplica por un factor \(a\),

\[
E'=\frac{E}{a^2}.
\]

Por tanto,

\[
a=2\Rightarrow E'=\frac{E}{4},
\]

\[
a=3\Rightarrow E'=\frac{E}{9},
\]

y para

\[
a=\frac12,
\]

se obtiene

\[
E'=\frac{E}{(1/2)^2}=4E.
\]

\end{explicacion}


\begin{ejercicio}

A una distancia de \(\SI{0.250}{\metre}\) de una carga puntual se mide un campo eléctrico de magnitud

\[
E=\SI{3.60e5}{\newton\per\coulomb}.
\]

El campo apunta alejándose de la carga.

Determine la carga fuente.

\end{ejercicio}

\begin{respuesta}

\[
\boxed{Q\approx+\SI{2.50}{\micro\coulomb}}
\]

\end{respuesta}

\begin{explicacion}

Partimos de

\[
E=k\frac{|Q|}{r^2}.
\]

Despejamos:

\[
|Q|=\frac{Er^2}{k}.
\]

Sustituyendo,

\[
|Q|=
\frac{(3.60\times10^5)(0.250)^2}{8.99\times10^9}
\approx\SI{2.50e-6}{\coulomb}.
\]

Por tanto,

\[
|Q|\approx\SI{2.50}{\micro\coulomb}.
\]

Como el campo apunta alejándose de la carga, la carga fuente debe ser positiva:

\[
\boxed{Q\approx+\SI{2.50}{\micro\coulomb}}.
\]

\end{explicacion}


\begin{ejercicio}

Indique el sentido del campo eléctrico producido por una carga puntual en cada caso:

\begin{enumerate}
    \item carga positiva;
    \item carga negativa.
\end{enumerate}

\end{ejercicio}

\begin{respuesta}

\begin{enumerate}
    \item Se aleja de la carga positiva.
    \item Apunta hacia la carga negativa.
\end{enumerate}

\end{respuesta}

\begin{explicacion}

La dirección del campo se define como la dirección de la fuerza que experimentaría una carga de prueba positiva.

Una carga fuente positiva repele a una carga de prueba positiva, por lo que el campo apunta hacia afuera.

Una carga fuente negativa atrae a una carga de prueba positiva, por lo que el campo apunta hacia la carga.

\end{explicacion}


\begin{ejercicio}

Dos cargas idénticas

\[
q_1=q_2=+\SI{4.00}{\micro\coulomb}
\]

se encuentran en

\[
x=-\SI{0.200}{\metre}
\]

y

\[
x=+\SI{0.200}{\metre}.
\]

Determine el campo eléctrico en el origen.

\end{ejercicio}

\begin{respuesta}

\[
\boxed{E=0}
\]

\end{respuesta}

\begin{explicacion}

Las dos cargas son positivas e idénticas y están a la misma distancia del origen.

Cada una produce un campo de igual magnitud:

\[
E_1=E_2.
\]

La carga situada a la izquierda produce un campo hacia la derecha, mientras que la situada a la derecha produce un campo hacia la izquierda.

Por tanto,

\[
E_x=E_1-E_2=0.
\]

Así,

\[
\boxed{E=0}.
\]

Cada carga produce campo en el origen, pero las dos contribuciones se cancelan.

\end{explicacion}


\begin{ejercicio}

Una carga

\[
Q=+\SI{2.00}{\micro\coulomb}
\]

se encuentra en el origen.

Determine el campo eléctrico en el punto

\[
P=(\SI{0.500}{\metre},0).
\]

\end{ejercicio}

\begin{respuesta}

\[
\boxed{E_x=+\SI{7.19e4}{\newton\per\coulomb}}
\]

\end{respuesta}

\begin{explicacion}

La distancia entre la carga y el punto es

\[
r=\SI{0.500}{\metre}.
\]

La magnitud del campo es

\[
E=k\frac{|Q|}{r^2}
=(8.99\times10^9)\frac{2.00\times10^{-6}}{(0.500)^2}.
\]

Entonces,

\[
E=\SI{7.19e4}{\newton\per\coulomb}.
\]

Como la carga es positiva y \(P\) está a su derecha, el campo apunta hacia \(+x\).

Por tanto,

\[
\boxed{E_x=+\SI{7.19e4}{\newton\per\coulomb}}.
\]

\end{explicacion}


\subsection{Nivel intermedio}

\begin{ejercicio}

Una carga \(q_1=+\SI{3.00}{\micro\coulomb}\) se encuentra en el origen y una carga
\(q_2=-\SI{4.00}{\micro\coulomb}\) se encuentra en
\((\SI{0.400}{\metre},0)\).

Determine el campo eléctrico en el punto

\[
P=(0,\SI{0.300}{\metre}).
\]

\end{ejercicio}

\begin{respuesta}

\[
\boxed{E_x\approx+\SI{1.15e5}{\newton\per\coulomb}}
\]

\[
\boxed{E_y\approx+\SI{2.13e5}{\newton\per\coulomb}}
\]

\[
\boxed{E\approx\SI{2.42e5}{\newton\per\coulomb}}
\]

\[
\boxed{\theta\approx61.7^\circ\text{ sobre }+x}
\]

\end{respuesta}

\begin{explicacion}

El campo de \(q_1\) apunta hacia arriba porque \(q_1\) es positiva.

La distancia entre \(q_1\) y \(P\) es

\[
r_1=\SI{0.300}{\metre}.
\]

Entonces,

\[
E_1=(8.99\times10^9)\frac{3.00\times10^{-6}}{(0.300)^2}
\approx\SI{3.00e5}{\newton\per\coulomb}.
\]

Sus componentes son

\[
E_{1x}=0,\qquad E_{1y}=+\SI{3.00e5}{\newton\per\coulomb}.
\]

Para \(q_2\),

\[
r_2=\sqrt{(0.400)^2+(0.300)^2}=\SI{0.500}{\metre}.
\]

Como \(q_2\) es negativa, el campo en \(P\) apunta hacia \(q_2\).

Su magnitud es

\[
E_2=(8.99\times10^9)\frac{4.00\times10^{-6}}{(0.500)^2}
\approx\SI{1.44e5}{\newton\per\coulomb}.
\]

De la geometría,

\[
\cos\theta=\frac{0.400}{0.500}=0.800,\qquad
\sin\theta=\frac{0.300}{0.500}=0.600.
\]

Por tanto,

\[
E_{2x}=+(1.44\times10^5)(0.800)\approx+\SI{1.15e5}{\newton\per\coulomb},
\]

\[
E_{2y}=-(1.44\times10^5)(0.600)\approx-\SI{8.63e4}{\newton\per\coulomb}.
\]

Aplicamos superposición:

\[
E_x=E_{1x}+E_{2x}\approx+\SI{1.15e5}{\newton\per\coulomb},
\]

\[
E_y=E_{1y}+E_{2y}\approx+\SI{2.13e5}{\newton\per\coulomb}.
\]

La magnitud es

\[
E=\sqrt{E_x^2+E_y^2}\approx\SI{2.42e5}{\newton\per\coulomb}.
\]

Como ambas componentes son positivas, el campo está en el primer cuadrante.

\[
\theta=\tan^{-1}\left(\frac{E_y}{E_x}\right)\approx61.7^\circ.
\]

\end{explicacion}


\begin{ejercicio}

Una carga

\[
q_1=+\SI{2.00}{\micro\coulomb}
\]

se encuentra en \((-\SI{0.300}{\metre},0)\), y una carga

\[
q_2=-\SI{5.00}{\micro\coulomb}
\]

se encuentra en \((0,\SI{0.400}{\metre})\).

Determine el campo eléctrico en el origen.

\end{ejercicio}

\begin{respuesta}

\[
\boxed{E_x\approx+\SI{2.00e5}{\newton\per\coulomb}}
\]

\[
\boxed{E_y\approx+\SI{2.81e5}{\newton\per\coulomb}}
\]

\[
\boxed{E\approx\SI{3.45e5}{\newton\per\coulomb}}
\]

\[
\boxed{\theta\approx54.6^\circ}
\]

\end{respuesta}

\begin{explicacion}

La carga \(q_1\) es positiva y está a la izquierda del origen, por lo que su campo en el origen apunta hacia \(+x\).

\[
E_1=(8.99\times10^9)\frac{2.00\times10^{-6}}{(0.300)^2}
\approx\SI{2.00e5}{\newton\per\coulomb}.
\]

Por tanto,

\[
E_{1x}=+\SI{2.00e5}{\newton\per\coulomb},\qquad E_{1y}=0.
\]

La carga \(q_2\) es negativa y está arriba del origen, por lo que el campo apunta hacia ella, es decir, hacia \(+y\).

\[
E_2=(8.99\times10^9)\frac{5.00\times10^{-6}}{(0.400)^2}
\approx\SI{2.81e5}{\newton\per\coulomb}.
\]

Así,

\[
E_{2x}=0,\qquad E_{2y}=+\SI{2.81e5}{\newton\per\coulomb}.
\]

Entonces,

\[
E_x=+\SI{2.00e5}{\newton\per\coulomb},
\qquad
E_y=+\SI{2.81e5}{\newton\per\coulomb}.
\]

La magnitud es

\[
E=\sqrt{(2.00\times10^5)^2+(2.81\times10^5)^2}
\approx\SI{3.45e5}{\newton\per\coulomb}.
\]

Finalmente,

\[
\theta=\tan^{-1}\left(\frac{2.81}{2.00}\right)\approx54.6^\circ.
\]

\end{explicacion}


\begin{ejercicio}

Una carga

\[
q_1=+\SI{3.00}{\micro\coulomb}
\]

se encuentra en \(x=-\SI{0.200}{\metre}\), y una carga

\[
q_2=-\SI{2.00}{\micro\coulomb}
\]

se encuentra en \(x=+\SI{0.300}{\metre}\).

Determine el campo eléctrico en el origen.

\end{ejercicio}

\begin{respuesta}

\[
\boxed{E_x\approx+\SI{8.74e5}{\newton\per\coulomb}}
\]

\end{respuesta}

\begin{explicacion}

La carga positiva situada a la izquierda produce un campo hacia la derecha:

\[
E_1=(8.99\times10^9)\frac{3.00\times10^{-6}}{(0.200)^2}
\approx\SI{6.74e5}{\newton\per\coulomb}.
\]

La carga negativa situada a la derecha también produce un campo hacia la derecha, porque el campo apunta hacia una carga negativa:

\[
E_2=(8.99\times10^9)\frac{2.00\times10^{-6}}{(0.300)^2}
\approx\SI{2.00e5}{\newton\per\coulomb}.
\]

Por tanto,

\[
E_x=E_1+E_2
\approx6.74\times10^5+2.00\times10^5.
\]

Así,

\[
\boxed{E_x\approx+\SI{8.74e5}{\newton\per\coulomb}}.
\]

\end{explicacion}


\begin{ejercicio}

Dos cargas positivas

\[
q_1=+\SI{4.00}{\micro\coulomb},
\qquad
q_2=+\SI{9.00}{\micro\coulomb}
\]

se encuentran en \(x=0\) y \(x=\SI{1.00}{\metre}\), respectivamente.

Determine el punto entre ellas donde el campo eléctrico es nulo.

\end{ejercicio}

\begin{respuesta}

\[
\boxed{x=\SI{0.400}{\metre}}
\]

medidos desde \(q_1\).

\end{respuesta}

\begin{explicacion}

Entre las cargas los campos tienen sentidos opuestos, por lo que pueden cancelarse.

Sea \(x\) la distancia desde \(q_1\).

La distancia hasta \(q_2\) es

\[
1.00-x.
\]

Para que el campo sea nulo,

\[
E_1=E_2.
\]

Entonces,

\[
k\frac{4.00}{x^2}
=
k\frac{9.00}{(1.00-x)^2}.
\]

Cancelamos \(k\):

\[
\frac{4}{x^2}=\frac{9}{(1-x)^2}.
\]

Tomando raíces positivas,

\[
\frac{2}{x}=\frac{3}{1-x}.
\]

Por tanto,

\[
2(1-x)=3x,
\]

\[
2=5x,
\]

y

\[
\boxed{x=\SI{0.400}{\metre}}.
\]

El punto queda más cerca de la carga de menor magnitud.

\end{explicacion}


\begin{ejercicio}

Dos cargas

\[
q_1=+\SI{4.00}{\micro\coulomb},
\qquad
q_2=-\SI{9.00}{\micro\coulomb}
\]

se encuentran en \(x=0\) y \(x=\SI{1.00}{\metre}\).

Determine el punto sobre el eje \(x\) donde el campo eléctrico es nulo.

\end{ejercicio}

\begin{respuesta}

\[
\boxed{x=-\SI{2.00}{\metre}}
\]

\end{respuesta}

\begin{explicacion}

Entre las cargas, ambos campos apuntan hacia la derecha, por lo que no pueden cancelarse.

El punto debe buscarse fuera del segmento, del lado de la carga de menor magnitud.

Sea \(a\) la distancia a la izquierda de \(q_1\). Entonces,

\[
r_1=a,\qquad r_2=1+a.
\]

La condición de cancelación es

\[
k\frac{4}{a^2}=k\frac{9}{(1+a)^2}.
\]

Simplificando,

\[
\frac{2}{a}=\frac{3}{1+a}.
\]

Entonces,

\[
2(1+a)=3a,
\]

de donde

\[
a=2.
\]

Como el punto está a la izquierda del origen,

\[
\boxed{x=-\SI{2.00}{\metre}}.
\]

\end{explicacion}


\begin{ejercicio}

En un punto del espacio, el campo eléctrico tiene componentes

\[
E_x=+\SI{1.20e5}{\newton\per\coulomb},
\qquad
E_y=-\SI{1.60e5}{\newton\per\coulomb}.
\]

Se coloca allí una carga

\[
q=-\SI{3.00}{\nano\coulomb}.
\]

Determine las componentes, magnitud y dirección de la fuerza eléctrica.

\end{ejercicio}

\begin{respuesta}

\[
\boxed{F_x=-\SI{3.60e-4}{\newton}},
\qquad
\boxed{F_y=+\SI{4.80e-4}{\newton}}
\]

\[
\boxed{F=\SI{6.00e-4}{\newton}}
\]

\[
\boxed{\theta\approx126.9^\circ\text{ desde }+x}
\]

\end{respuesta}

\begin{explicacion}

Trabajamos directamente por componentes:

\[
F_x=qE_x,\qquad F_y=qE_y.
\]

Entonces,

\[
F_x=(-3.00\times10^{-9})(1.20\times10^5)
=-\SI{3.60e-4}{\newton},
\]

\[
F_y=(-3.00\times10^{-9})(-1.60\times10^5)
=+\SI{4.80e-4}{\newton}.
\]

La magnitud es

\[
F=\sqrt{F_x^2+F_y^2}.
\]

Así,

\[
F=\sqrt{(3.60\times10^{-4})^2+(4.80\times10^{-4})^2}
=\SI{6.00e-4}{\newton}.
\]

Como

\[
F_x<0,\qquad F_y>0,
\]

la fuerza está en el segundo cuadrante.

El ángulo de referencia es

\[
\theta_r=\tan^{-1}\left(\frac{4.80}{3.60}\right)\approx53.1^\circ.
\]

Por tanto,

\[
\theta=180^\circ-53.1^\circ=126.9^\circ.
\]

\end{explicacion}


\begin{ejercicio}

Dos cargas iguales

\[
q_1=q_2=+\SI{2.00}{\micro\coulomb}
\]

se encuentran en

\[
(0,+\SI{0.300}{\metre})
\]

y

\[
(0,-\SI{0.300}{\metre}).
\]

Determine el campo eléctrico en

\[
P=(\SI{0.400}{\metre},0).
\]

\end{ejercicio}

\begin{respuesta}

\[
\boxed{E_x\approx+\SI{1.15e5}{\newton\per\coulomb}},
\qquad
\boxed{E_y=0}
\]

\end{respuesta}

\begin{explicacion}

La distancia desde cualquiera de las cargas hasta \(P\) es

\[
r=\sqrt{(0.400)^2+(0.300)^2}=\SI{0.500}{\metre}.
\]

Cada carga produce un campo de magnitud

\[
E_0=(8.99\times10^9)\frac{2.00\times10^{-6}}{(0.500)^2}
=\SI{7.19e4}{\newton\per\coulomb}.
\]

De la geometría,

\[
\cos\theta=\frac{0.400}{0.500}=0.800,
\qquad
\sin\theta=\frac{0.300}{0.500}=0.600.
\]

Las componentes verticales tienen igual magnitud y sentidos opuestos:

\[
E_y=E_0(0.600)-E_0(0.600)=0.
\]

Las componentes horizontales apuntan ambas hacia la derecha:

\[
E_x=2E_0(0.800).
\]

Entonces,

\[
E_x=2(7.19\times10^4)(0.800)
\approx\SI{1.15e5}{\newton\per\coulomb}.
\]

Por tanto,

\[
\boxed{E_x\approx+\SI{1.15e5}{\newton\per\coulomb}},
\qquad
\boxed{E_y=0}.
\]

\end{explicacion}


\begin{ejercicio}

Una carga

\[
q_1=+\SI{2.00}{\micro\coulomb}
\]

se encuentra en \(x=-\SI{0.200}{\metre}\).

¿Qué carga \(q_2\), situada en \(x=+\SI{0.400}{\metre}\), hará que el campo eléctrico en el origen sea cero?

\end{ejercicio}

\begin{respuesta}

\[
\boxed{q_2=+\SI{8.00}{\micro\coulomb}}
\]

\end{respuesta}

\begin{explicacion}

La carga \(q_1\) es positiva y produce en el origen un campo hacia la derecha:

\[
E_1=k\frac{2.00\times10^{-6}}{(0.200)^2}.
\]

Para cancelar este campo, \(q_2\), situada a la derecha, debe producir un campo hacia la izquierda.

Por tanto, \(q_2\) también debe ser positiva.

La condición es

\[
E_1=E_2.
\]

Entonces,

\[
k\frac{2.00\times10^{-6}}{(0.200)^2}
=
k\frac{q_2}{(0.400)^2}.
\]

Cancelamos \(k\) y despejamos:

\[
q_2
=
(2.00\times10^{-6})
\frac{(0.400)^2}{(0.200)^2}.
\]

Así,

\[
q_2=8.00\times10^{-6}\,\si{\coulomb}.
\]

Por tanto,

\[
\boxed{q_2=+\SI{8.00}{\micro\coulomb}}.
\]

\end{explicacion}


\begin{ejercicio}

Un anillo de radio

\[
a=\SI{0.200}{\metre}
\]

posee una carga uniforme

\[
Q=+\SI{6.00}{\micro\coulomb}.
\]

Determine el campo eléctrico en un punto de su eje situado a

\[
x=\SI{0.150}{\metre}
\]

del centro.

\end{ejercicio}

\begin{respuesta}

\[
\boxed{E_x\approx+\SI{5.18e5}{\newton\per\coulomb}}
\]

\end{respuesta}

\begin{explicacion}

El campo sobre el eje de un anillo uniformemente cargado es

\[
E_x=k\frac{Qx}{(x^2+a^2)^{3/2}}.
\]

Sustituyendo,

\[
E_x=
(8.99\times10^9)
\frac{(6.00\times10^{-6})(0.150)}
{\left[(0.150)^2+(0.200)^2\right]^{3/2}}.
\]

Como

\[
(0.150)^2+(0.200)^2=0.0625,
\]

se obtiene

\[
\boxed{E_x\approx+\SI{5.18e5}{\newton\per\coulomb}}.
\]

El signo positivo corresponde al lado \(x>0\) de un anillo con carga positiva.

\end{explicacion}


\begin{ejercicio}

Una varilla de longitud

\[
L=\SI{0.500}{\metre}
\]

contiene una carga uniforme total

\[
Q=+\SI{4.00}{\micro\coulomb}.
\]

Determine:

\begin{enumerate}
    \item la densidad lineal de carga;
    \item la carga contenida en un segmento de longitud \(\SI{2.00}{\centi\metre}\).
\end{enumerate}

\end{ejercicio}

\begin{respuesta}

\[
\boxed{\lambda=\SI{8.00}{\micro\coulomb\per\metre}}
\]

\[
\boxed{\Delta q=\SI{0.160}{\micro\coulomb}}
\]

\end{respuesta}

\begin{explicacion}

Para una distribución uniforme,

\[
\lambda=\frac{Q}{L}.
\]

Entonces,

\[
\lambda=\frac{4.00\times10^{-6}}{0.500}
=\SI{8.00e-6}{\coulomb\per\metre}.
\]

Por tanto,

\[
\boxed{\lambda=\SI{8.00}{\micro\coulomb\per\metre}}.
\]

El segmento tiene longitud

\[
\Delta\ell=\SI{0.0200}{\metre}.
\]

Como

\[
\Delta q=\lambda\Delta\ell,
\]

resulta

\[
\Delta q=(8.00\times10^{-6})(0.0200)
=\SI{1.60e-7}{\coulomb}.
\]

Por tanto,

\[
\boxed{\Delta q=\SI{0.160}{\micro\coulomb}}.
\]

\end{explicacion}


\subsection{Nivel avanzado}

\begin{ejercicio}

Tres cargas ocupan los vértices de un triángulo equilátero de lado

\[
a=\SI{0.300}{\metre}.
\]

En el vértice inferior izquierdo se encuentra

\[
q_1=+\SI{3.00}{\micro\coulomb},
\]

en el inferior derecho,

\[
q_2=-\SI{6.00}{\micro\coulomb}.
\]

Determine el campo eléctrico en el vértice superior.

\end{ejercicio}

\begin{respuesta}

\[
\boxed{E_x\approx+\SI{4.50e5}{\newton\per\coulomb}}
\]

\[
\boxed{E_y\approx-\SI{2.60e5}{\newton\per\coulomb}}
\]

\[
\boxed{E\approx\SI{5.19e5}{\newton\per\coulomb}}
\]

\[
\boxed{\theta=30.0^\circ\text{ por debajo de }+x}
\]

\end{respuesta}

\begin{explicacion}

Ambas cargas están a una distancia

\[
r=\SI{0.300}{\metre}
\]

del vértice superior.

El campo producido por \(q_1\) tiene magnitud

\[
E_1=(8.99\times10^9)\frac{3.00\times10^{-6}}{(0.300)^2}
\approx\SI{3.00e5}{\newton\per\coulomb}.
\]

Como \(q_1\) es positiva, el campo se aleja de ella y forma \(60^\circ\) sobre \(+x\):

\[
E_{1x}=E_1\cos60^\circ,
\qquad
E_{1y}=E_1\sin60^\circ.
\]

Para \(q_2\),

\[
E_2=(8.99\times10^9)\frac{6.00\times10^{-6}}{(0.300)^2}
\approx\SI{5.99e5}{\newton\per\coulomb}.
\]

Como \(q_2\) es negativa, su campo apunta hacia ella, es decir, \(60^\circ\) por debajo de \(+x\):

\[
E_{2x}=E_2\cos60^\circ,
\qquad
E_{2y}=-E_2\sin60^\circ.
\]

Entonces,

\[
E_x=(3.00\times10^5)(0.500)+(5.99\times10^5)(0.500)
\approx\SI{4.50e5}{\newton\per\coulomb},
\]

\[
E_y=(3.00\times10^5)(0.866)-(5.99\times10^5)(0.866)
\approx-\SI{2.60e5}{\newton\per\coulomb}.
\]

La magnitud es

\[
E=\sqrt{E_x^2+E_y^2}\approx\SI{5.19e5}{\newton\per\coulomb}.
\]

El campo está en el cuarto cuadrante y

\[
\theta_r=\tan^{-1}\left(\frac{|E_y|}{E_x}\right)=30.0^\circ.
\]

Por tanto,

\[
\boxed{\theta=30.0^\circ\text{ por debajo de }+x}.
\]

\end{explicacion}


\begin{ejercicio}

Cuatro cargas ocupan los vértices de un cuadrado de lado

\[
a=\SI{0.400}{\metre}.
\]

Tres de ellas tienen carga

\[
+\SI{2.00}{\micro\coulomb},
\]

mientras que la carga del vértice superior derecho es

\[
-\SI{2.00}{\micro\coulomb}.
\]

Determine el campo eléctrico en el centro del cuadrado.

\end{ejercicio}

\begin{respuesta}

\[
\boxed{E_x\approx+\SI{3.18e5}{\newton\per\coulomb}}
\]

\[
\boxed{E_y\approx+\SI{3.18e5}{\newton\per\coulomb}}
\]

\[
\boxed{E\approx\SI{4.50e5}{\newton\per\coulomb}}
\]

La dirección es \(45^\circ\) sobre \(+x\).

\end{respuesta}

\begin{explicacion}

La distancia desde cualquier vértice al centro es

\[
r=\frac{a}{\sqrt2}
=\frac{0.400}{\sqrt2}
\approx\SI{0.283}{\metre}.
\]

Cada carga produce un campo de igual magnitud:

\[
E_0=k\frac{2.00\times10^{-6}}{r^2}
\approx\SI{2.25e5}{\newton\per\coulomb}.
\]

Cada campo forma \(45^\circ\) con los ejes, por lo que la magnitud de cada componente es

\[
E_c=E_0\cos45^\circ
\approx\SI{1.59e5}{\newton\per\coulomb}.
\]

Sumamos por componentes.

En \(x\),

\[
E_x=E_c-E_c+E_c+E_c=2E_c.
\]

Así,

\[
E_x\approx+\SI{3.18e5}{\newton\per\coulomb}.
\]

En \(y\),

\[
E_y=E_c+E_c-E_c+E_c=2E_c,
\]

por lo que

\[
E_y\approx+\SI{3.18e5}{\newton\per\coulomb}.
\]

La magnitud es

\[
E=\sqrt{E_x^2+E_y^2}
\approx\SI{4.50e5}{\newton\per\coulomb}.
\]

Como ambas componentes son iguales y positivas,

\[
\boxed{\theta=45^\circ}.
\]

El campo apunta hacia el vértice que contiene la carga negativa.

\end{explicacion}


\begin{ejercicio}

Dos cargas

\[
q_1=q_2=+\SI{2.00}{\micro\coulomb}
\]

se encuentran en

\[
(\SI{0.300}{\metre},+\SI{0.400}{\metre})
\]

y

\[
(\SI{0.300}{\metre},-\SI{0.400}{\metre}).
\]

Una tercera carga \(q_3\) se coloca en

\[
(-\SI{0.200}{\metre},0).
\]

Determine \(q_3\) para que el campo eléctrico en el origen sea cero.

\end{ejercicio}

\begin{respuesta}

\[
\boxed{q_3\approx+\SI{0.384}{\micro\coulomb}}
\]

\end{respuesta}

\begin{explicacion}

Las cargas \(q_1\) y \(q_2\) están a

\[
r=\sqrt{(0.300)^2+(0.400)^2}=\SI{0.500}{\metre}
\]

del origen.

Cada una produce un campo de magnitud

\[
E_0=(8.99\times10^9)\frac{2.00\times10^{-6}}{(0.500)^2}
=\SI{7.19e4}{\newton\per\coulomb}.
\]

Las componentes verticales se cancelan por simetría.

La componente horizontal de cada campo es

\[
E_{0x}=-E_0\frac{0.300}{0.500}
=-0.600E_0.
\]

Por tanto,

\[
E_{12,x}=-2(0.600)(7.19\times10^4)
\approx-\SI{8.63e4}{\newton\per\coulomb}.
\]

Para cancelar este campo, \(q_3\) debe producir

\[
E_{3x}=+\SI{8.63e4}{\newton\per\coulomb}.
\]

Como \(q_3\) está a la izquierda del origen, debe ser positiva para producir un campo hacia la derecha.

Aplicamos

\[
E_3=k\frac{q_3}{(0.200)^2}.
\]

Despejando,

\[
q_3=\frac{E_3(0.200)^2}{k}.
\]

Entonces,

\[
q_3=
\frac{(8.63\times10^4)(0.200)^2}{8.99\times10^9}
\approx3.84\times10^{-7}\,\si{\coulomb}.
\]

Por tanto,

\[
\boxed{q_3\approx+\SI{0.384}{\micro\coulomb}}.
\]

\end{explicacion}


\begin{ejercicio}

Una varilla delgada de longitud \(L\) está centrada sobre el eje \(x\) y posee una densidad lineal uniforme positiva \(\lambda\).

Determine el campo eléctrico en el punto

\[
P=(0,a)
\]

situado sobre la mediatriz de la varilla.

\end{ejercicio}

\begin{respuesta}

\[
\boxed{E_x=0}
\]

\[
\boxed{
E_y=
k\frac{\lambda L}
{a\sqrt{a^2+(L/2)^2}}
}
\]

\end{respuesta}

\begin{explicacion}

Tomamos un elemento de longitud \(dx\) ubicado en la posición \(x\).

Su carga es

\[
dq=\lambda\,dx.
\]

La distancia desde el elemento hasta \(P\) es

\[
r=\sqrt{x^2+a^2}.
\]

El campo infinitesimal tiene magnitud

\[
dE=k\frac{dq}{r^2}
=k\frac{\lambda\,dx}{x^2+a^2}.
\]

Por simetría, para cada elemento situado en \(+x\) existe otro en \(-x\).

Sus componentes horizontales se cancelan:

\[
\boxed{E_x=0}.
\]

Las componentes verticales se suman.

Como

\[
\cos\theta=\frac{a}{\sqrt{x^2+a^2}},
\]

la componente vertical es

\[
dE_y=dE\frac{a}{\sqrt{x^2+a^2}}.
\]

Entonces,

\[
dE_y=
k\lambda a
\frac{dx}{(x^2+a^2)^{3/2}}.
\]

Integramos sobre toda la varilla:

\[
E_y=
k\lambda a
\int_{-L/2}^{L/2}
\frac{dx}{(x^2+a^2)^{3/2}}.
\]

Por simetría,

\[
E_y=
2k\lambda a
\int_0^{L/2}
\frac{dx}{(x^2+a^2)^{3/2}}.
\]

Utilizando

\[
\int\frac{dx}{(x^2+a^2)^{3/2}}
=
\frac{x}{a^2\sqrt{x^2+a^2}},
\]

obtenemos

\[
E_y=
2k\lambda a
\left[
\frac{x}{a^2\sqrt{x^2+a^2}}
\right]_0^{L/2}.
\]

Simplificando,

\[
\boxed{
E_y=
k\frac{\lambda L}
{a\sqrt{a^2+(L/2)^2}}
}.
\]

\end{explicacion}


\begin{ejercicio}

Una varilla uniformemente cargada tiene

\[
L=\SI{0.400}{\metre},
\qquad
\lambda=+\SI{3.00}{\micro\coulomb\per\metre}.
\]

Está centrada sobre el eje \(x\).

Determine el campo eléctrico en el punto

\[
P=(0,\SI{0.300}{\metre}).
\]

\end{ejercicio}

\begin{respuesta}

\[
\boxed{E_x=0}
\]

\[
\boxed{E_y\approx+\SI{9.97e4}{\newton\per\coulomb}}
\]

\end{respuesta}

\begin{explicacion}

Por simetría,

\[
E_x=0.
\]

Para una varilla uniforme,

\[
E_y=
k\frac{\lambda L}
{a\sqrt{a^2+(L/2)^2}}.
\]

Sustituimos

\[
a=\SI{0.300}{\metre},
\qquad
L=\SI{0.400}{\metre},
\qquad
\lambda=\SI{3.00e-6}{\coulomb\per\metre}.
\]

Entonces,

\[
E_y=
(8.99\times10^9)
\frac{(3.00\times10^{-6})(0.400)}
{(0.300)\sqrt{(0.300)^2+(0.200)^2}}.
\]

Por tanto,

\[
\boxed{E_y\approx+\SI{9.97e4}{\newton\per\coulomb}}.
\]

El signo positivo corresponde a una varilla positiva y un punto situado sobre ella en \(y>0\).

\end{explicacion}


\begin{ejercicio}

Una varilla uniformemente cargada ocupa el intervalo

\[
0\leq x\leq L
\]

sobre el eje \(x\), con densidad lineal positiva \(\lambda\).

Determine el campo eléctrico en un punto situado sobre el mismo eje, a una distancia \(d\) del extremo derecho de la varilla.

\end{ejercicio}

\begin{respuesta}

\[
\boxed{
E_x=k\lambda\left(\frac{1}{d}-\frac{1}{L+d}\right)
}
\]

\[
\boxed{E_y=0}
\]

\end{respuesta}

\begin{explicacion}

El punto de observación se encuentra en

\[
x=L+d.
\]

Consideramos un elemento \(dq=\lambda\,dx\) situado en una posición \(x\).

La distancia al punto es

\[
r=L+d-x.
\]

Como la carga es positiva y el punto está a la derecha de toda la varilla, todas las contribuciones apuntan hacia \(+x\).

Por tanto,

\[
dE_x=k\frac{\lambda\,dx}{(L+d-x)^2}.
\]

Integramos desde \(x=0\) hasta \(x=L\):

\[
E_x=
k\lambda
\int_0^L
\frac{dx}{(L+d-x)^2}.
\]

El resultado es

\[
E_x=
k\lambda
\left(
\frac{1}{d}
-
\frac{1}{L+d}
\right).
\]

Por tanto,

\[
\boxed{
E_x=k\lambda\left(\frac{1}{d}-\frac{1}{L+d}\right)
}.
\]

Como todas las contribuciones están sobre el eje \(x\),

\[
\boxed{E_y=0}.
\]

\end{explicacion}


\begin{ejercicio}

Una semicircunferencia de radio

\[
R=\SI{0.200}{\metre}
\]

posee una carga uniforme total

\[
Q=+\SI{4.00}{\micro\coulomb}.
\]

La semicircunferencia se encuentra sobre la mitad superior del plano \(xy\) y tiene su centro en el origen.

Determine el campo eléctrico en el centro.

\end{ejercicio}

\begin{respuesta}

\[
\boxed{E_x=0}
\]

\[
\boxed{E_y\approx-\SI{5.72e5}{\newton\per\coulomb}}
\]

\end{respuesta}

\begin{explicacion}

La densidad lineal es

\[
\lambda=\frac{Q}{\pi R}.
\]

Tomamos un elemento de arco

\[
d\ell=R\,d\theta,
\]

por lo que

\[
dq=\lambda R\,d\theta.
\]

La distancia de todos los elementos al centro es \(R\), así que

\[
dE=k\frac{dq}{R^2}.
\]

Por simetría, las componentes horizontales se cancelan:

\[
\boxed{E_x=0}.
\]

Las componentes verticales apuntan hacia abajo.

Por tanto,

\[
dE_y=-dE\sin\theta.
\]

Sustituyendo,

\[
dE_y
=
-\frac{k\lambda}{R}\sin\theta\,d\theta.
\]

Integramos desde \(0\) hasta \(\pi\):

\[
E_y
=
-\frac{k\lambda}{R}
\int_0^\pi\sin\theta\,d\theta.
\]

Como

\[
\int_0^\pi\sin\theta\,d\theta=2,
\]

resulta

\[
E_y=-\frac{2k\lambda}{R}.
\]

Usando

\[
\lambda=\frac{Q}{\pi R},
\]

obtenemos

\[
E_y=-\frac{2kQ}{\pi R^2}.
\]

Sustituyendo,

\[
E_y=
-\frac{2(8.99\times10^9)(4.00\times10^{-6})}
{\pi(0.200)^2}.
\]

Finalmente,

\[
\boxed{E_y\approx-\SI{5.72e5}{\newton\per\coulomb}}.
\]

\end{explicacion}


\begin{ejercicio}

Un anillo de radio

\[
a=\SI{0.200}{\metre}
\]

posee una carga

\[
Q=+\SI{6.00}{\micro\coulomb}.
\]

Compare el campo exacto sobre su eje a una distancia

\[
x=\SI{2.00}{\metre}
\]

con el campo que se obtendría tratando toda la carga como una carga puntual en el centro.

\end{ejercicio}

\begin{respuesta}

Campo exacto:

\[
\boxed{E_{\text{anillo}}\approx\SI{1.33e4}{\newton\per\coulomb}}
\]

Aproximación puntual:

\[
\boxed{E_{\text{puntual}}\approx\SI{1.35e4}{\newton\per\coulomb}}
\]

La diferencia relativa es aproximadamente

\[
\boxed{1.50\%}.
\]

\end{respuesta}

\begin{explicacion}

El campo exacto del anillo es

\[
E_{\text{anillo}}
=
k\frac{Qx}{(x^2+a^2)^{3/2}}.
\]

Sustituyendo,

\[
E_{\text{anillo}}
=
(8.99\times10^9)
\frac{(6.00\times10^{-6})(2.00)}
{\left[(2.00)^2+(0.200)^2\right]^{3/2}}.
\]

Entonces,

\[
E_{\text{anillo}}\approx\SI{1.33e4}{\newton\per\coulomb}.
\]

Si tratamos el anillo como una carga puntual,

\[
E_{\text{puntual}}=k\frac{Q}{x^2}.
\]

Por tanto,

\[
E_{\text{puntual}}
=
(8.99\times10^9)
\frac{6.00\times10^{-6}}{(2.00)^2}
\approx\SI{1.35e4}{\newton\per\coulomb}.
\]

La diferencia relativa respecto del valor exacto es

\[
\frac{E_{\text{puntual}}-E_{\text{anillo}}}
{E_{\text{anillo}}}\times100\%
\approx1.50\%.
\]

Como

\[
x=10a,
\]

el punto se encuentra suficientemente lejos para que el anillo se comporte aproximadamente como una carga puntual.

\end{explicacion}


\begin{ejercicio}

Una distribución lineal de carga ocupa el segmento del eje \(x\) comprendido entre \(x=a\) y \(x=b\), con

\[
0<a<b.
\]

Su densidad lineal no es uniforme y viene dada por

\[
\lambda(x)=\alpha x,
\]

donde \(\alpha>0\).

Determine el campo eléctrico en el origen.

\end{ejercicio}

\begin{respuesta}

\[
\boxed{
E_x=-k\alpha\ln\left(\frac{b}{a}\right)
}
\]

\[
\boxed{E_y=0}
\]

\end{respuesta}

\begin{explicacion}

Un elemento de longitud \(dx\) tiene carga

\[
dq=\lambda(x)\,dx=\alpha x\,dx.
\]

La distancia del elemento al origen es

\[
r=x.
\]

Como toda la distribución es positiva y se encuentra a la derecha del origen, cada elemento produce en el origen un campo hacia \(-x\).

Por tanto,

\[
dE_x=-k\frac{dq}{x^2}.
\]

Sustituyendo \(dq\),

\[
dE_x=-k\frac{\alpha x\,dx}{x^2}
=-k\alpha\frac{dx}{x}.
\]

Integramos desde \(a\) hasta \(b\):

\[
E_x=-k\alpha\int_a^b\frac{dx}{x}.
\]

Así,

\[
E_x=-k\alpha[\ln x]_a^b.
\]

Por tanto,

\[
\boxed{
E_x=-k\alpha\ln\left(\frac{b}{a}\right)
}.
\]

Como toda la distribución está sobre el eje \(x\),

\[
\boxed{E_y=0}.
\]

\end{explicacion}


\begin{ejercicio}

Una semicircunferencia superior de radio

\[
R=\SI{0.200}{\metre}
\]

posee una carga uniforme

\[
Q=+\SI{4.00}{\micro\coulomb}.
\]

Además, una carga puntual positiva \(q\) se coloca sobre el eje \(y\) en

\[
y=-\SI{0.300}{\metre}.
\]

Determine \(q\) para que el campo eléctrico total en el origen sea cero.

\end{ejercicio}

\begin{respuesta}

\[
\boxed{q\approx+\SI{5.73}{\micro\coulomb}}
\]

\end{respuesta}

\begin{explicacion}

El campo de la semicircunferencia en el centro tiene únicamente componente vertical.

Para una semicircunferencia positiva situada sobre el origen,

\[
E_{\text{arco},y}
=
-\frac{2kQ}{\pi R^2}.
\]

Sustituyendo,

\[
E_{\text{arco},y}
=
-\frac{2(8.99\times10^9)(4.00\times10^{-6})}
{\pi(0.200)^2}
\approx-\SI{5.72e5}{\newton\per\coulomb}.
\]

Para que el campo total sea cero, la carga puntual debe producir en el origen un campo hacia arriba de igual magnitud.

Como la carga está debajo del origen, debe ser positiva para que su campo apunte alejándose de ella, es decir, hacia \(+y\).

Entonces,

\[
E_q=k\frac{q}{(0.300)^2}.
\]

Imponemos

\[
E_q=5.72\times10^5.
\]

Despejando,

\[
q=
\frac{(5.72\times10^5)(0.300)^2}{8.99\times10^9}.
\]

Por tanto,

\[
\boxed{q\approx+\SI{5.73}{\micro\coulomb}}.
\]

Así,

\[
E_y=E_{q,y}+E_{\text{arco},y}=0.
\]

Por simetría,

\[
E_x=0.
\]

\end{explicacion}